%\clearpage

\section {Introduction}
\label{sec:intro}
In this technical note, we report the results of a 3-flavour $\nu/\bar{\nu}$ joint analysis performed by the {\em VALOR} group, on the combined \RunPeriod dataset, using data taken in both neutrino-enhanced forward horn current (FHC) mode (\TotalPotNu{} POT), and antineutrino-enhanced reverse horn current (RHC) mode (\TotalPotNuBar{} POT).

% TODO: Write about the changes since the last analysis, incuding the xsec model, POT, osc param priors, flux tuning etc.
%M: I was not sure if  there should be any detailed description in this section. Hope this version is fine 

Compared to the previous analysis, random phase approximation  (BeRPA) was replaced with 8 Q2 parameters. Instead of having one CCOther shape parameter we moved to CCDIS_BY(based on Bodek-Yang model), CCMpi_BY, CCmisc, CCDIS nu and nubar normalizations. Also, a Coulomb shift correction was introduced for calculation of the Ereco for CCQE events. Implemented 4 (2 for nu and 2 for nubar) energy dependent 2p2h parameters as well as 2(nu and nubar) CC normalization parameters.  This analysis includes new nuclear model. FSI parameters were moved  from SK detector covariance matrix to BANFF matrix  All these changes are further described in the section \cref {sec:systematics}
As for the flux tuning, the flux version used in the current analysis is based on the 2009 replica target tuning with the improvement method  of treating run-by-run correlations of p beam profile systematic.Since the last last analysis oscillation priors for the \ssqthonethree and \ssqthonetwo were changed according to the PDG19  release. 



% TODO: Update this for run 10; Are we updating the MH study?

Details of the Super-K samples used and the analysis procedure can be found in \cite{OurRun18JointAnalysis3Flavour}, \cref{sec:spectra} presents predicted spectra, \cref{sec:systematics} describes the systematic errors considered and their effects on the extrapolated single-ring event spectra, whilst the sensitivity of the various fits is considered in \cref{sec:sensitivity} before presenting the fit results in \cref{sec:results}, including the Feldman-Cousins confidence intervals for $\delta_{CP}$, a comparison of the fit results to the previous run in \cref{sec:comparison}, \color{blue} a study on the frequentist properties of posterior probability for mass ordering is presented in \cref{sec:mh}\color{black} and a summary is given in \cref{sec:summary}.

\clearpage
